\subsection*{Phase 2 haplotypes, atlas, and 2La analysis}

The mosquito atlas used the open Ag1000G Phase 2 AR1 phased-haplotype release on AgamP4 \citep{malariagen_phase2_ar1_2017,ag1000g2020}. We prespecified nine established species-by-population cohorts with at least 40 samples. Populations GM, GW, and KE were excluded from the primary atlas because their species status was reported as uncertain. To keep the panel identical across chromosome arms, eligible samples were restricted to the 1,058 identifiers present in all five released HDF5 files. Forty diploids per cohort were selected without replacement using NumPy seed 2026. Complete, biallelic, segregating sites were retained, and the fixed broadened-\(N_e\) checkpoint was applied without retraining. Adjacent-SNP rates were aggregated into 50- and 100-kb windows.

The 2La interval was 20.524--42.166 Mb. Tag-SNP dosage was averaged and rounded following the published MalariaGEN karyotyping rule. Arrangement frequency used every eligible released sample in each population, whereas map inference used the frozen 40-diploid panel. Suppression depth was \(1-r_{\mathrm{in}}/r_{\mathrm{out}}\), using median positive 100-kb rates. Two-sided Pearson and Spearman tests related suppression depth to \(2p(1-p)\).

\subsection*{Phase 2 cross comparison}

All 11 Phase 2 crosses and autosomes 2R, 2L, 3R, and 3L were analyzed from the released SHAPEIT haplotypes. Sites informative for one parent required that parent to be heterozygous, the mate homozygous, at least eight offspring callable, and a Mendelian-error fraction no greater than 0.05. SHAPEIT values above one (codes 2 or 255) were treated as uncalled. At most one marker per 10-kb bin was retained. Adjacent markers were oriented from sibling transmissions, and a new phase block was started when more than one sibling supported the less common orientation. A two-state Viterbi caller used genotype-error probability 0.0025, a fixed 1-cM/Mb transition prior, and at least three markers in each flanking state run. Breakpoint intervals wider than 1 Mb were excluded.

Direct rates divided event counts by detection-aware parental-transmission exposure. Windows required at least 80\% first-to-last marker-span coverage, and both direct and inferred recombination maps were divided by their chromosome-arm means before a pooled 5-Mb Spearman comparison. Because the Phase 2 release lacks the later genotype-quality and explicit site-filter arrays, this analysis is interpreted only as a broad release-specific comparison.

\subsection*{Resistance-region controls}

The 15 AgamP4 resistance regions were defined from phenotype markers and named, uniquely mappable mechanism genes in the literature. Candidate controls were on the same chromosome arm, at least 0.50 Mb from any focal region, and matched by empirical-percentile distance in telomere position, log SNP density, nucleotide diversity, and H12. H12 used at most 128 deterministically spaced SNPs in each local window. Up to 75 controls were retained per locus. The population statistic was the exponent of the mean log focal-to-control ratio, and 5,000 within-population draws provided descriptive permutation values. All summaries used the same frozen set of 15 resistance regions. Species-stratified intervals in Fig.~\ref{fig:anopheles}F were obtained from 4,000 deterministic population-bootstrap samples and describe variation among cohorts rather than uncertainty across independently replicated species.
